Four linked models organize the formal argument, in place of a single universal optimization equation. They are kept distinct because each can fail independently of the others.

\section*{Admissible Resource Contribution}

Each node $i$ has a set $\Aset_i(t)$ of permitted actions, formed as the intersection of an owner policy and a fixed safety envelope as developed in Chapter~\ref{ch:idle}, so that it reflects local constraints, safety requirements, ownership, and primary function, and every action taken must satisfy
\[
a_i(t)\in\Aset_i(t).
\]

\section*{Verified Value}

A task produces value only when its result satisfies a specified validation procedure. Numerical throughput alone is insufficient.

\section*{Lifetime Resource Accounting}

Incremental energy, capital costs, hardware replacement, emissions, and the value of useful outputs are compared over a specified horizon, as in the discounted sum of Chapter~\ref{ch:horizon}.

\section*{Metric Corruption Under Optimization Pressure}

The fourth model examines how rewards based on measurable proxies can cause participants to expand resource consumption or manipulate reported performance. Proposition~\ref{prop:goodhart} is its simplest instance.

\medskip

A node may be admissible but economically unproductive. A computation may be scientifically valuable but too expensive. An efficient operation may violate local authority. A seemingly successful reward system may be optimizing a corrupted metric. A strong treatment resists reducing these distinctions to a single scalar objective.
